% Options for packages loaded elsewhere
% Options for packages loaded elsewhere
\PassOptionsToPackage{unicode}{hyperref}
\PassOptionsToPackage{hyphens}{url}
\PassOptionsToPackage{dvipsnames,svgnames,x11names}{xcolor}
%
\documentclass[
  12pt,
  letterpaper,
  DIV=11,
  numbers=noendperiod]{scrreprt}
\usepackage{xcolor}
\usepackage[margin=1in]{geometry}
\usepackage{amsmath,amssymb}
\setcounter{secnumdepth}{5}
\usepackage{iftex}
\ifPDFTeX
  \usepackage[T1]{fontenc}
  \usepackage[utf8]{inputenc}
  \usepackage{textcomp} % provide euro and other symbols
\else % if luatex or xetex
  \usepackage{unicode-math} % this also loads fontspec
  \defaultfontfeatures{Scale=MatchLowercase}
  \defaultfontfeatures[\rmfamily]{Ligatures=TeX,Scale=1}
\fi
\usepackage{lmodern}
\ifPDFTeX\else
  % xetex/luatex font selection
  \setmainfont[]{Arial}
\fi
% Use upquote if available, for straight quotes in verbatim environments
\IfFileExists{upquote.sty}{\usepackage{upquote}}{}
\IfFileExists{microtype.sty}{% use microtype if available
  \usepackage[]{microtype}
  \UseMicrotypeSet[protrusion]{basicmath} % disable protrusion for tt fonts
}{}
\makeatletter
\@ifundefined{KOMAClassName}{% if non-KOMA class
  \IfFileExists{parskip.sty}{%
    \usepackage{parskip}
  }{% else
    \setlength{\parindent}{0pt}
    \setlength{\parskip}{6pt plus 2pt minus 1pt}}
}{% if KOMA class
  \KOMAoptions{parskip=half}}
\makeatother
% Make \paragraph and \subparagraph free-standing
\makeatletter
\ifx\paragraph\undefined\else
  \let\oldparagraph\paragraph
  \renewcommand{\paragraph}{
    \@ifstar
      \xxxParagraphStar
      \xxxParagraphNoStar
  }
  \newcommand{\xxxParagraphStar}[1]{\oldparagraph*{#1}\mbox{}}
  \newcommand{\xxxParagraphNoStar}[1]{\oldparagraph{#1}\mbox{}}
\fi
\ifx\subparagraph\undefined\else
  \let\oldsubparagraph\subparagraph
  \renewcommand{\subparagraph}{
    \@ifstar
      \xxxSubParagraphStar
      \xxxSubParagraphNoStar
  }
  \newcommand{\xxxSubParagraphStar}[1]{\oldsubparagraph*{#1}\mbox{}}
  \newcommand{\xxxSubParagraphNoStar}[1]{\oldsubparagraph{#1}\mbox{}}
\fi
\makeatother


\usepackage{longtable,booktabs,array}
\newcounter{none} % for unnumbered tables
\usepackage{calc} % for calculating minipage widths
% Correct order of tables after \paragraph or \subparagraph
\usepackage{etoolbox}
\makeatletter
\patchcmd\longtable{\par}{\if@noskipsec\mbox{}\fi\par}{}{}
\makeatother
% Allow footnotes in longtable head/foot
\IfFileExists{footnotehyper.sty}{\usepackage{footnotehyper}}{\usepackage{footnote}}
\makesavenoteenv{longtable}
\usepackage{graphicx}
\makeatletter
\newsavebox\pandoc@box
\newcommand*\pandocbounded[1]{% scales image to fit in text height/width
  \sbox\pandoc@box{#1}%
  \Gscale@div\@tempa{\textheight}{\dimexpr\ht\pandoc@box+\dp\pandoc@box\relax}%
  \Gscale@div\@tempb{\linewidth}{\wd\pandoc@box}%
  \ifdim\@tempb\p@<\@tempa\p@\let\@tempa\@tempb\fi% select the smaller of both
  \ifdim\@tempa\p@<\p@\scalebox{\@tempa}{\usebox\pandoc@box}%
  \else\usebox{\pandoc@box}%
  \fi%
}
% Set default figure placement to htbp
\def\fps@figure{htbp}
\makeatother


% definitions for citeproc citations
\NewDocumentCommand\citeproctext{}{}
\NewDocumentCommand\citeproc{mm}{%
  \begingroup\def\citeproctext{#2}\cite{#1}\endgroup}
\makeatletter
 % allow citations to break across lines
 \let\@cite@ofmt\@firstofone
 % avoid brackets around text for \cite:
 \def\@biblabel#1{}
 \def\@cite#1#2{{#1\if@tempswa , #2\fi}}
\makeatother
\newlength{\cslhangindent}
\setlength{\cslhangindent}{1.5em}
\newlength{\csllabelwidth}
\setlength{\csllabelwidth}{3em}
\newenvironment{CSLReferences}[2] % #1 hanging-indent, #2 entry-spacing
 {\begin{list}{}{%
  \setlength{\itemindent}{0pt}
  \setlength{\leftmargin}{0pt}
  \setlength{\parsep}{0pt}
  % turn on hanging indent if param 1 is 1
  \ifodd #1
   \setlength{\leftmargin}{\cslhangindent}
   \setlength{\itemindent}{-1\cslhangindent}
  \fi
  % set entry spacing
  \setlength{\itemsep}{#2\baselineskip}}}
 {\end{list}}
\usepackage{calc}
\newcommand{\CSLBlock}[1]{\hfill\break\parbox[t]{\linewidth}{\strut\ignorespaces#1\strut}}
\newcommand{\CSLLeftMargin}[1]{\parbox[t]{\csllabelwidth}{\strut#1\strut}}
\newcommand{\CSLRightInline}[1]{\parbox[t]{\linewidth - \csllabelwidth}{\strut#1\strut}}
\newcommand{\CSLIndent}[1]{\hspace{\cslhangindent}#1}



\setlength{\emergencystretch}{3em} % prevent overfull lines

\providecommand{\tightlist}{%
  \setlength{\itemsep}{0pt}\setlength{\parskip}{0pt}}



 


\usepackage{booktabs}
\usepackage{longtable}
\usepackage{array}
\usepackage{multirow}
\usepackage{wrapfig}
\usepackage{float}
\usepackage{colortbl}
\usepackage{pdflscape}
\usepackage{tabu}
\usepackage{threeparttable}
\usepackage{threeparttablex}
\usepackage[normalem]{ulem}
\usepackage{makecell}
\usepackage{xcolor}
\KOMAoption{captions}{tablesignature}
\makeatletter
\@ifpackageloaded{tcolorbox}{}{\usepackage[skins,breakable]{tcolorbox}}
\@ifpackageloaded{fontawesome5}{}{\usepackage{fontawesome5}}
\definecolor{quarto-callout-color}{HTML}{909090}
\definecolor{quarto-callout-note-color}{HTML}{0758E5}
\definecolor{quarto-callout-important-color}{HTML}{CC1914}
\definecolor{quarto-callout-warning-color}{HTML}{EB9113}
\definecolor{quarto-callout-tip-color}{HTML}{00A047}
\definecolor{quarto-callout-caution-color}{HTML}{FC5300}
\definecolor{quarto-callout-color-frame}{HTML}{acacac}
\definecolor{quarto-callout-note-color-frame}{HTML}{4582ec}
\definecolor{quarto-callout-important-color-frame}{HTML}{d9534f}
\definecolor{quarto-callout-warning-color-frame}{HTML}{f0ad4e}
\definecolor{quarto-callout-tip-color-frame}{HTML}{02b875}
\definecolor{quarto-callout-caution-color-frame}{HTML}{fd7e14}
\makeatother
\makeatletter
\@ifpackageloaded{bookmark}{}{\usepackage{bookmark}}
\makeatother
\makeatletter
\@ifpackageloaded{caption}{}{\usepackage{caption}}
\AtBeginDocument{%
\ifdefined\contentsname
  \renewcommand*\contentsname{Table of contents}
\else
  \newcommand\contentsname{Table of contents}
\fi
\ifdefined\listfigurename
  \renewcommand*\listfigurename{List of Figures}
\else
  \newcommand\listfigurename{List of Figures}
\fi
\ifdefined\listtablename
  \renewcommand*\listtablename{List of Tables}
\else
  \newcommand\listtablename{List of Tables}
\fi
\ifdefined\figurename
  \renewcommand*\figurename{Figure}
\else
  \newcommand\figurename{Figure}
\fi
\ifdefined\tablename
  \renewcommand*\tablename{Table}
\else
  \newcommand\tablename{Table}
\fi
}
\@ifpackageloaded{float}{}{\usepackage{float}}
\floatstyle{ruled}
\@ifundefined{c@chapter}{\newfloat{codelisting}{h}{lop}}{\newfloat{codelisting}{h}{lop}[chapter]}
\floatname{codelisting}{Listing}
\newcommand*\listoflistings{\listof{codelisting}{List of Listings}}
\makeatother
\makeatletter
\makeatother
\makeatletter
\@ifpackageloaded{caption}{}{\usepackage{caption}}
\@ifpackageloaded{subcaption}{}{\usepackage{subcaption}}
\makeatother
\usepackage{bookmark}
\IfFileExists{xurl.sty}{\usepackage{xurl}}{} % add URL line breaks if available
\urlstyle{same}
\hypersetup{
  pdftitle={Study Title: A Content Analysis of Cozy Games and Gender Performance},
  pdfauthor={Gerri Marion},
  colorlinks=true,
  linkcolor={darkred},
  filecolor={Maroon},
  citecolor={Blue},
  urlcolor={Blue},
  pdfcreator={LaTeX via pandoc}}


\title{Study Title: A Content Analysis of Cozy Games and Gender
Performance}
\author{Gerri Marion}
\date{2026-05-01}
\begin{document}
\maketitle

\renewcommand*\contentsname{Table of contents}
{
\setcounter{tocdepth}{2}
\tableofcontents
}

\bookmarksetup{startatroot}

\chapter{Abstract}\label{abstract}

\textbf{Author:} Gerri Marion \textbf{Instructor:} Dr.~Alex P. Leith
\textbf{Institution:} Southern Illinois University Edwardsville
\textbf{Date:} May 2026

\begin{center}\rule{0.5\linewidth}{0.5pt}\end{center}

\textbf{Abstract}

\emph{This study examines how YouTube gameplay videos covered the cozy
game genre during 04/01/2026-04/05/2026. Using a two-part content
analysis (N = 100), this study first employed qualitative analysis to
identify emergent themes in the coverage, then tested relationships
among coded variables using the Chi-Square test. Findings indicate that
there was no statistical significange between gender and gaming
mechanics, nor through the qualitative's results of gender performance
and gaming mechanics. These results suggest that cozy games are a
relaxing experience that allows a release of the pressures of gender
performance, encouring the inclusive space for everyone. It suggests
that because of the lack of pressures and complications of cozy games,
there is not much of a reason to perform by their gender identity.}

\begin{center}\rule{0.5\linewidth}{0.5pt}\end{center}

\subsection{About This Study}\label{about-this-study}

This project presents original content analysis research conducted in
partial fulfillment of MC 451: Research Methods in Mass Media at SIUE. I
looked into the cozy game genre and how gender performance is impacted
through a genre marketed towards women. This study was conducted using
YouTube gameplay videos of \emph{Stardew Valley} and \emph{Fields of
Mistria}.

\begin{center}\rule{0.5\linewidth}{0.5pt}\end{center}

\subsection{How to Navigate This
Report}\label{how-to-navigate-this-report}

{\def\LTcaptype{none} % do not increment counter
\begin{longtable}[]{@{}
  >{\raggedright\arraybackslash}p{(\linewidth - 2\tabcolsep) * \real{0.5000}}
  >{\raggedright\arraybackslash}p{(\linewidth - 2\tabcolsep) * \real{0.5000}}@{}}
\toprule\noalign{}
\begin{minipage}[b]{\linewidth}\raggedright
Section
\end{minipage} & \begin{minipage}[b]{\linewidth}\raggedright
What You Will Find
\end{minipage} \\
\midrule\noalign{}
\endhead
\bottomrule\noalign{}
\endlastfoot
\href{01-literature-review.qmd}{Literature Review} & Background, theory,
and research questions \\
\href{02-methodology.qmd}{Methodology} & Data source, sampling,
codebook, and procedures \\
\href{03-study-1.qmd}{Study 1: Qualitative Analysis} & Emergent themes,
patterns, and variable development \\
\href{04-study-2.qmd}{Study 2: Quantitative Analysis} & Descriptive
statistics, visualization, and inferential testing \\
\href{05-discussion.qmd}{Discussion} & Interpretation, limitations, and
conclusion \\
\href{references.qmd}{References} & Full bibliography \\
\end{longtable}
}

\bookmarksetup{startatroot}

\chapter{Literature Review}\label{literature-review}

The cozy game genre has been around for decades, targeting their cute
aesthetics and ``simpler'' game play to women and minority communities.
It gives them the space to truly express themselves and be relieved of
the pressure and harassment from online gaming communities, specifically
from men. However, there is very little studying how Butler's (1988)
Gender Performativity Theory plays into a genre known for it's soft and
simple game play. This study looks to see if gender is still performed
when new players turn to the cozy game genre, addressing it through a
qualitative and quantitative content analysis.

\section{Background and Significance}\label{background-and-significance}

Cozy games have been known in the video game community for years, with
games such as Electronic Arts \emph{The Sims} series and Amccus
\emph{Harvest Moon}, now known as \emph{Story of Seasons}. The biggest
reason why this genre had a rapid incline is how women and minorities
were and still are being treated in online gaming communities. The most
notable of specifically women being targed in online gaming spaces, the
majority being by men, is the GamerGate scandal in 2014. According to
Cote (2020), GamerGate started as ``{[}s{]}purred by (unfounded)
accusations that game designer Zoe Quinn had exchanged sexual favors for
positive coverage of her game \emph{Depression Quest}, the online
Twitter campaign supposedly advocated for better ethics in game
journalism. In practice, it directed a slew of harassment at female game
developers, game studies scholars, journalists, players, and allies.''
These women were harassed, received death threats, doxxed, and so on.
Later, a study was conducted by Tang \& Fox (2016) looking at how men
interact with women and minorities in online gaming spaces. They found
that ``certain personality favors and contextual factors predict men's
harassment behavior in online video games.'' These included topics such
as intense sexual and general harassment towards those individuals.

With recurring themes of harassment and harmful targeting towards women
and other minorities, the cozy game genre rose in popularity with these
target demographics. The COVID-19 pandemic was the final push for cozy
games to gain momentum as ``a surge of women joining the hobby and
vocalized support for games that do not align with what has been deemed
as mainstream.'' Brtis (2022)

While there are plenty of studies that look into why the cozy game genre
exists and why it's so popular, there's very little that looks into how
gender is performed when these kinds of games are played. A study
conducted by {``Beyond the {`Core-Gamer'}''} (2015) focused on the
gender preferences of different gaming genres and found that men heavily
leaned towards action and shooter style of games, or what's known as
``core-gamer'' kind of games, while older people and women preferred
puzzle and card games. Fang (2025) studied gender performance, but it
wasn't through the cozy game genre. They used \emph{Honor of Kings}
(TiMi Studio Group, 2015), a popular online game in China that's mostly
known for being ``a male-dominated virtual space'' and wanted to see how
different genders performed in that area. They found themes of
gender-switching, ``gendered strategies,'' and different forms of gender
performance depending on the gender they chose for their character.
``Women employ `gender-neutral' self representation, competitive
character choices, skill-focused gaming to demonstrate the agency and
power inherent to gender performativity'' while men can simply play the
game as themselves, ``enjoy the novelty of role-reversal, experiencing
the game from the female perspective.'' Fang (2025) Even though Fang
found some themes of gender performance, it was in an online game known
for players who identify as men. There isn't much that addresses gender
performance from a genre that encase players who identify as women or
other gender identities and backgrounds. The goal for this study,
especially after Fang's findings, was to turn the tables and see if
different genders preferred specific mechanics and how they acted out
their genders through a game viewed as cozy.

One thing to note is that \emph{Honor of Kings} could be considered as a
cozy game for those who aren't men and play it as there is no official
agreement between players, game designers, journalists, etc., that can
properly define what counts as a cozy game; it's subjective based on the
person and their experiences. Boudreau et al. (2025) However, for the
sake of this study and being stated in Fang's (2025) article that
\emph{Honor of Kings} is mostly for masculine players, cozy games will
be focused more on the targeted demographic of women and minorities.

The chosen media for this study was looking at YouTube gameplay videos
about a player's first experiences in cozy games. Concerned Ape's 2016
title \emph{Stardew Valley} not only is the game people in the community
find as the top cozy game, but ``imagines sexual and cultural
progressiveness as a long-term part of traditional and rural society.''
Mackay \& Roberts (2025) The other game chosen was NPC Studio's
\emph{Fields of Mistria} (2024), a game with very similar mechanics in a
90's anime style, also inspired by games like \emph{Stardew Valley}.

Both of these games are very similar in mechanics, romance options,
inclusivity and theme, making them the ideals for what cozy games are
usually thought of. By looking at YouTube gameplay videos of first
experiences with these kinds of games, we'll be able to see if their
gender is performed throughout the video and how. This looks at what
gaming mechanics a player leans towards, how they interact and verbally
respond to livestream chats and characters, if they spend time a lot of
time in character creation and farm decoration, and how their gender is
performed, if at all.

\section{Theoretical Framework}\label{theoretical-framework}

This study is grounded in Butler (1988) Gender Performativity Theory.
The key idea is that gender is a performance, repeated by the society
and culture that individual is raised in. If their performance is done
correctly or incorrectly, they are praised or punished by those who
``accurately'' perform their gender identity. According to Butler, the
repeated actions of that performance ``is a reenactment and
reexperiencing of a set of meanings already socially established; it is
the mundane and ritualized form of of their legitimation.'' (p.~526)

Applied to the present study, Gender Performativity Theory would
continue through a genre that already has a predominantly feminine
demographic. This study predicts that women and minorities will be more
relaxed and willing to embrace their femininity through creativity, such
as the amount of time spent creating a character and farm
exterior-interior decorating, while men will lean more towards hands-on
gaming mechanics, such as mining and combat and fishing, and spend less
time on their farm. They will perform their genders based on repeated
actions praised in their personal lives, and possibly through audience
interaction if the video is a part of a livestream.

The article examined in the Annotated Manuscript --- Boudreau et al.
(2025) --- provides a semi-direct application of this framework. While
Boudreau focuses on Social Construction Theory when it came to her dive
into cozy games and what they are, the idea of gender performance is
socially constructed. Butler confirms this through the \emph{repeated}
actions of a person, something that would exist within the construction
of that society. With cozy games being predominantly women and
minorities, it allows and encourages that idea of gender performance
through what makes the person feel ``cozy,'' though again, the idea of
``cozy'' is based on personal experiences.

\section{Research Questions}\label{research-questions}

Based on the gaps in the existing literature and the logic of Gender
Performativity Theory, this study addresses the following research
questions:

\begin{quote}
\textbf{RQ1:} How would YouTube gameplay videos frame how different
genders play cozy games?
\end{quote}

\begin{quote}
\textbf{RQ2:} What is the relationship between gender performance and
how the interact with their first ``100 Days in \emph{Stardew Valley} or
\emph{Fields of Mistria}'' YouTube gameplay videos?
\end{quote}

Together, these questions aim to produce a systematic, empirically
grounded account of how YouTube gameplay videos covers the cozy game
genre --- one that moves beyond anecdotal observation toward replicable
measurement.

\bookmarksetup{startatroot}

\chapter{Methodology}\label{methodology}

This study employed a two-part content analysis to address the research
questions introduced in the previous section. Study 1 used qualitative
methods --- immersion, memo writing, and thematic analysis --- to
identify emergent patterns in the source material and develop the coding
variables used in Study 2. Study 2 applied those variables
systematically through quantitative content analysis, producing
descriptive statistics and inferential tests. Content analysis is
appropriate here because it provides a replicable, transparent method
for identifying patterns in media content at a scale not feasible
through close reading alone (\textbf{Neuendorf2017?};
\textbf{Krippendorff2018?}).

\section{Data Source}\label{data-source}

Data was collected from YouTube gameplay videos of a player's first
experience in \emph{Stardew Valley} and \emph{Fields of Mistria}. This
source was selected because my main focus of this study is how gender
and gender performance plays into the cozy game genre, if at all. The
genre itself was created during and after COVID as a response to toxic
masculinity in online gaming communities. By choosing these two games,
it allowed the ability to see how new players interact with the cozy
game genre as well as being able to see if some of those gender
performances would reveal during gameplay. The YouTube gameplay videos
of both \emph{Stardew Valley} and \emph{Fields of Mistria} ``first
times'' provides access to typical gaming mechanics in cozy games, the
possibility of gender being performed through them, and being recent
enough to have accessible data, making it well suited to the systematic
sampling procedure described below.

\section{Sample and Sampling
Procedure}\label{sample-and-sampling-procedure}

Items were identified using the search string
\texttt{("stardew\ valley"\ OR\ "fields\ of\ mistria")\ AND\ ("first\ playthrough"\ OR\ "blind\ playthrough"\ OR\ "first\ time")},
applied to YouTube videos with results restricted to 01/01/2016 through
04/01/2026. This date range was selected because it involved the release
of \emph{Stardew Valley} and \emph{Fields of Mistria} while also still
being able to reveal gender performance through in-game mechanics.
Future updates did not impact the overall first interactions of a
person's gameplay.

The initial search returned 152,000 items:

\begin{itemize}
\tightlist
\item
  ``All items meeting the inclusion criteria were retained, producing a
  final sample of 100 items.''
\end{itemize}

Items were excluded if they were non-commentary videos, videos where the
player did not speak for the entire video, or only spoke a few words
before remaining silent for the rest of the video.

The final sample consisted of \textbf{N = 100} items, each constituting
one unit of analysis: One YouTube gameplay video, as defined by Google's
Support website as ``videos that contain in-game footage to show a good
sense of how a game is played.''

\section{Analytical Approach}\label{analytical-approach}

The analysis proceeded in two phases. In Study 1, a qualitative
immersion process was used to identify emergent themes, develop
conceptual categories, and construct operational definitions for each
variable. This inductive process --- grounded in the memo-writing and
immersion techniques described in the course textbook --- ensured that
the coding scheme reflected patterns present in the data rather than
being imposed from the outside.

In Study 2, the resulting codebook was applied systematically to the
full sample. Each item was coded on 5 variables, producing a structured
dataset that was analyzed using descriptive statistics and an
inferential test appropriate to the research questions (described in
Study 2).

\section{Variables and Measures}\label{variables-and-measures}

\textbf{Gender} refers to ``The characteristics of women, men, girls and
boys that are socially constructed. This include norms, behaviors and
roles associated with being a woman, man, girl or boy, as well as
relationships with each other.'' \emph{Gender and Health} (n.d.) For the
purposes of this study, Gender was operationalized as viewing the
beginning of the video, up until after the character creation, and
selecting one of two values. Coders assigned one of two values: 1 =
Masculine, 2 = Feminine. If a player chose something different from
their gender identity or showed discomfort in selective gender choices,
the decision rule is to code 0 = gender identity conflict. Items for
which a determination could not be made were assigned code 9 and
excluded from analysis.

\textbf{Gaming Mechanics} refers to ``The specific rules, actions, and
systems that define how players interact with a game. Determines what a
player can do, how the game responds, and how process, challenge, and
rewards are created.'' {admin} (2026). This variable was operationalized
as watching the last fifteen minutes \emph{Stardew Valley} or
\emph{Fields of Mistria} gameplay video, or the ``100th day'' of their
video entirely, Winter 16. Then, noting how many times a day one
category is chosen, selecting the highest total number's category at the
end. This will be their ``dominating'' gaming mechanic., coded into 5
categories: 1 = Land, 2 = Mining, 3 = Community, 4 = Romance, if
applicable, 5 = Fishing. This variable had a few decision rules, such as
a code 4 being a manual write-down of what NPC was chosen for marriage,
upgrading farm tools is a code 1 while processing geodes is a code 2,
and anything that had equal tallies in multiple categories were coded as
a 9 and looked into.

\textbf{Gameplay Videos} refers to ``Videos that contain in-game footage
to show a good sense of how a game is played.'' \emph{About Gameplay
Videos - {Google Ads Help}} (n.d.) This variable was operationalized
Coded based on watching the last fifteen minutes of a \emph{Stardew
Valley} or \emph{Fields of Mistria} gameplay video, or the ``100th day''
of their video entirely, or Winter 16, then and choosing one of four
categories: 1 = Full Video Without Help, 2 = Full Video With Help, 3 =
Playlist Without Help, 4 = Playlist With Help. If a player pauses for a
period of time, such as a streamer stopping to talk to their chat or a
player pausing to take a break, skip through it until they begin to play
again.

\textbf{Gender Performance} refers to gender as not only an identity,
but repeated, social performance using actions that, when done correctly
or incorrectly, are praised or punished by those same actors a part of
that performance. Butler (1988) This variable was operationalized as
based on what gaming mechanics players lean towards based on their
gender identity and performance of their gender and assigning one of two
values, coded into 2 categories: 1 = Yes, they did perform by their
gender, 2 = No, they did not perform by their gender. If the player
seemed indifferent towards gender identity and performance and instead
played the game neutrally or for enjoyment, code = 3.

\textbf{Farm Choice} refers to the various different farm map options a
player has at the beginning of the game. There are eight to choose from
in the current version of \emph{Stardew Valley}, however as \emph{Fields
of Mistria} is still in early access, there is currently only one farm
type in the game. This variable was operationalized as at the beginning
of the video during character customization. This is an option off to
the side that a player can choose from, then coded into 2 categories: 1
= Standard Farm, 2 = Other, or a different farm choice option. If the
game was \emph{Fields of Mistria}, the code = 0. This is a game still in
early access and, at the time of writing this, does not have options for
alternative farm layouts.

\textbf{Season} refers to the seasonal time period the player is in
during their most recent video. This variable was operationalized as at
the start of the data collection time, coded into 4 categories: 1 =
Spring, 2 = Summer, 3 = Fall or Autumn. 4 = Winter. If unable to tell at
the start or throughout the video and if a calendar was absent for icon
symbol of season, code = 9.

\section{Pilot Test and Codebook
Refinement}\label{pilot-test-and-codebook-refinement}

Prior to full data collection, the codebook was applied to a pilot
sample of 10 items to identify ambiguities and refine decision rules.
The pilot revealed that multiple variables needed an update or to be
redefined. Gameplay Videos showed that not every video could be watched
the whole way through, leading to the singular day or last fifteen
minutes of their video. Gaming Mechanics did not originally have a
section for fishing or how to handle shops and geode processing, so
those were updated along with category 5 creation. Gender Performance
also took into account verbal responses and interactions with chat
livestream or responses to NPCs as part of the performance, so those
were revised to include it for masculine, feminine, and indifference.
Farm Choice did not consider how \emph{Fields of Mistria} has a standard
farm for everyone, so code = 0 was created for that category. Finally,
the Season variable was added and its four categories to see how far
each player made it in their gameplays. No further substantial revisions
were required, and the pilot items were retained in the final sample.

\bookmarksetup{startatroot}

\chapter{Study 1: Qualitative Content
Analysis}\label{study-1-qualitative-content-analysis}

The quantitative analysis presented in the next chapter depends on a
coding scheme --- a set of variables with clearly defined categories.
But those variables did not emerge from thin air. Before any systematic
coding could begin, a period of qualitative immersion was necessary to
understand the source material on its own terms and identify the
patterns worth measuring. This chapter documents that process.

\section{Immersion Process}\label{immersion-process}

To develop the coding scheme, all 100 items from the sample were read
closely and without a predetermined framework. The goal was to approach
the material inductively --- to notice what was there rather than impose
categories from the outside. During this process, observational memos
were written to capture initial reactions, recurring patterns, and
unexpected features of the coverage.

When immersing myself in every individual video, what stood out the most
is how gaming mechanics played into a player's in-game day-to-day, as
well as how far along each person played the two games. Most players
were in the winter season, specifically around my Winter 16/100th day
marker, and they had well-developed farms and goods overall. There was
also a notable amount of my code 9 when it came to gaming mechanics,
mostly because the different gaming mechanics would end up as a tie
during tallying and data collecting. A theme I noticed is that players
were either indifferent when it came to gender performance and played a
cozy game for the fact that it was fun and relaxing, or acted out their
gender's verbally. This was grouped into the gender performance category
and noted as much as possible.

\section{Emergent Themes}\label{emergent-themes}

Three primary themes emerged from the qualitative analysis:

\subsection{Theme 1: In-Game Season}\label{theme-1-in-game-season}

One of the variables that was the most curious is how far along would
every person be in their gameplay videos. This could be a factor when it
came to livestreams or playlist videos, encouraging content creation and
interest, but it became a strong theme. Winter 16 was the cut-off for
these videos as it would be a ``100th day'' for every player, and most
in the dataset did make it to winter. If the player wasn't in Winter,
the next dominant season was Spring with only a few videos or one single
livestream to have the full experience.

\subsection{Theme 2: Gaming Mechanics Code =
9}\label{theme-2-gaming-mechanics-code-9}

Gaming mechanics was one thing that needed to be looked into when it
came to the overall study. While ``community'' and ``land'' were
relatively high categories, mostly coming from the repetitive gameplay
of farmland tools and activities as well as quest completion and
gift-giving, there was a handful fo code = 9 areas. These were code =
9's because they ended up tying in multiple gameplay mechanics either
throughout a single in-game day or through a ``100 Days in/of'' YouTube
video. These code = 9's tied mostly in community and land, which are the
prominent mechanics for a player to do in these two gaming titles.

\subsection{Theme 3: Verbal Responses}\label{theme-3-verbal-responses}

This was not a variable considered in the study as gender performance
focused on the gaming mechanics and interactions of in-game NPC's and
not through a chat livestream or YouTube commentary video. However,
there were noted moments of language usually noticed in online gaming
communities throughout some of the videos. Instances such as men making
innuendos or focusing on ``the grind'' while women interacted with the
community the most, as well as prioritizing gift-giving to build in-game
romances. This theme could be explored further with future studies. This
theme tied into the gender performance variable as it did seem to bring
some praise from others as well as making the initial player feel good
about their act.

\section{From Themes to Variables}\label{from-themes-to-variables}

The themes identified above informed the construction of the coding
variables used in Study 2. Specifically:

\begin{itemize}
\item
  \textbf{In-Game Season} was operationalized as the variable
  \textbf{Season}, coded as 1 = Spring, 2 = Summer, 3 = Fall, 4 =
  Winter. The qualitative analysis revealed that season and longevity of
  a character's game file played a significant part. While the
  progression of each game could've been encouraged due to popularity
  for a player's channel, they still showed enough interest to keep the
  gameplay going. Winter ended up being noted the most with Spring
  coming in second. If the player was in Spring, the video was most
  likely about the first week and the player's first experiences
  overall.
\item
  \textbf{Gaming Mechanics Code = 9} was operationalized as the variable
  \textbf{Gaming Mechanics}, coded as 1 = Land, 2 = Mining, 3 =
  Community, 4 = Romance, 5 = Fishing, and 9 = Undetermined. 9 stayed as
  ``undetermined'' due to the fact that, at the end of data collection
  for a video, specific gaming mechanics were equal with land and
  community being the highest of those ties. This could've been it's own
  code number, but due to the fact that it happened infrequently, a code
  = 9 made the most sense. It also made sense in the overall since the
  main gaming mechanics of both \emph{Stardew Valley} and \emph{Fields
  of Mistria} is running a farm and interacting with the entire land of
  on-farm and out-of-farm maps and the interactions of NPC's.
\item
  \textbf{Verbal Responses} was operationalized as the variable
  \textbf{Gender Performance}, coded as 1 = did perform their gender, 2
  = did not perform their gender, 3 = indifferent about gender
  performance. Verbal responses were included and made note of as part
  of gender performance. Themes of how different genders interacted with
  in-game mechanics and NPC's as well as how they interacted with their
  own videos and chats showed some themes of gender performance.
\end{itemize}

This inductive process --- moving from open-ended reading to structured
observation --- ensured that the coding scheme reflected the actual
patterns present in the data rather than being imposed from a textbook
or prior study. The resulting variables formed the basis of the
systematic quantitative analysis presented in the next chapter.

\section{Edge Cases and Coding
Decisions}\label{edge-cases-and-coding-decisions}

Several items resisted easy classification. The variable
\textbf{Gameplay Videos} showed the greatest problem during the pilot
testing, then later throughout the rest of my coding data. Some players
had their videos as multiple livestreams in a single playlist, others
had full-length videos of ``100 Days in/of'' either gaming title. Some
required additional assistance, such as in-game chat or the Wikipedia
pages for both games, while others preferred the challenge and sought to
figure things out on their own. Because of the length of time for the
study and each video, it was noted how their videos were organized
through one of four categories. Then, the data collection went to the
most recent video and studying the last 15-minutes, or seeking out as
close to Winter 16 as possible. These edge cases and the decision rules
developed to handle them were documented in the codebook and applied
consistently throughout the quantitative coding phase.

\bookmarksetup{startatroot}

\chapter{Study 2: Quantitative Content
Analysis}\label{study-2-quantitative-content-analysis}

Using the variables developed through the qualitative analysis in Study
1, the full sample (N = 100 items) was coded systematically. This
chapter presents the quantitative results: descriptive statistics
characterizing the sample, followed by an inferential test addressing
the study's research questions.

\begin{center}\rule{0.5\linewidth}{0.5pt}\end{center}

\section{Descriptive Statistics}\label{descriptive-statistics}

\subsection{Overview of the Sample}\label{overview-of-the-sample}

\begingroup\fontsize{13}{15}\selectfont

\begin{longtable}[t]{lcc}

\toprule
\cellcolor[HTML]{b71c1c}{\textcolor{white}{\textbf{Source / Outlet}}} & \cellcolor[HTML]{b71c1c}{\textcolor{white}{\textbf{N}}} & \cellcolor[HTML]{b71c1c}{\textcolor{white}{\textbf{\% of Sample}}}\\
\midrule
did perform & 55 & 55\%\\
did not perform & 6 & 6\%\\
indifferent & 36 & 36\%\\
undetermined & 3 & 3\%\\
\bottomrule


\caption{\label{tbl-sample-overview}Table 1. Sample composition by
Gender Performance.}

\tabularnewline
\end{longtable}

\endgroup{}

The data came from YouTube gameplay videos, watching them the whole way
through, then deciding whether or not the different genders showed
gender performance. It was distributed through the span of five days
with fairly consistent video upload times.

\begin{center}\rule{0.5\linewidth}{0.5pt}\end{center}

\subsection{Distribution of Primary Variable: Gender
Performance}\label{distribution-of-primary-variable-gender-performance}

\begingroup\fontsize{13}{15}\selectfont

\begin{longtable}[t]{lccc}

\toprule
\cellcolor[HTML]{b71c1c}{\textcolor{white}{\textbf{Gender Performance}}} & \cellcolor[HTML]{b71c1c}{\textcolor{white}{\textbf{N}}} & \cellcolor[HTML]{b71c1c}{\textcolor{white}{\textbf{Percent}}} & \cellcolor[HTML]{b71c1c}{\textcolor{white}{\textbf{Cumulative \%}}}\\
\midrule
did perform & 55 & 55\% & 55\%\\
indifferent & 36 & 36\% & 91\%\\
did not perform & 6 & 6\% & 97\%\\
undetermined & 3 & 3\% & 100\%\\
\bottomrule


\caption{\label{tbl-var1-freq}Table 2. Frequency distribution of Gender
Performance (N = 100).}

\tabularnewline
\end{longtable}

\endgroup{}

The most common when it came to gender performance is that half of my
dataset did perform their gender, with signs of indifference being the
second most common. The least common is those who did not perform by
their gender. This connects to Study 1 when it came to the verbal
responses and how that coincided with gender performance.

\begin{center}\rule{0.5\linewidth}{0.5pt}\end{center}

\subsection{Visualization: Gender
Performance}\label{visualization-gender-performance}

\begin{figure}

\centering{

\pandocbounded{\includegraphics[keepaspectratio]{04-study-2_files/figure-pdf/fig-var1-bar-1.pdf}}

}

\caption{\label{fig-var1-bar}Figure 1. Distribution of Gender
Performance across all coded items. The key visual pattern is those who
did perform their gender and those who either did not or showed
indifference to gender identity}

\end{figure}%

The dominant category are those who did perform by their gender. The
distribution is fairly concentrated in one category and spread out in
the others. The most surprising is how little those did not perform by
their gender and were indifferent overall. They were more focused on
enjoying the game without gender constraints even though they had a
gender identity.

\begin{center}\rule{0.5\linewidth}{0.5pt}\end{center}

\subsection{Distribution of Secondary Variable: Gaming
Mechanics}\label{distribution-of-secondary-variable-gaming-mechanics}

\begingroup\fontsize{13}{15}\selectfont

\begin{longtable}[t]{lcc}

\toprule
\cellcolor[HTML]{b71c1c}{\textcolor{white}{\textbf{Gaming Mechanics}}} & \cellcolor[HTML]{b71c1c}{\textcolor{white}{\textbf{N}}} & \cellcolor[HTML]{b71c1c}{\textcolor{white}{\textbf{Percent}}}\\
\midrule
community & 40 & 40\%\\
land & 36 & 36\%\\
mining & 12 & 12\%\\
undetermined & 10 & 10\%\\
fishing & 2 & 2\%\\
\bottomrule


\caption{\label{tbl-var2-freq}Table 3. Frequency distribution of Gaming
Mechanics.}

\tabularnewline
\end{longtable}

\endgroup{}

The top gaming mechanic is the community gaming mechanic with land
following behind it. Fishing and romance were the lowest ones, to my
surprise. While there were interests in romance, it did end up being a
part of community or did not make it far enough video-wise to have an
official romantic interest. Community and land being in the top make
sense because both games prioritize those mechanics the most;
interacting with NPC's and working around your farm land or interacting
with the entire map.

\begin{center}\rule{0.5\linewidth}{0.5pt}\end{center}

\section{Inferential Analysis}\label{inferential-analysis}

\subsection{Chi-Square Test of Independence: Gender Performance × Gaming
Mechanic}\label{chi-square-test-of-independence-gender-performance-gaming-mechanic}

\begingroup\fontsize{13}{15}\selectfont

\begin{longtable}[t]{lcccclc}

\toprule
\multicolumn{1}{c}{ } & \multicolumn{6}{c}{Gaming Mechanic} \\
\cmidrule(l{3pt}r{3pt}){2-7}
\cellcolor[HTML]{b71c1c}{\textcolor{white}{\textbf{Gender Performance}}} & \cellcolor[HTML]{b71c1c}{\textcolor{white}{\textbf{land}}} & \cellcolor[HTML]{b71c1c}{\textcolor{white}{\textbf{mining}}} & \cellcolor[HTML]{b71c1c}{\textcolor{white}{\textbf{community}}} & \cellcolor[HTML]{b71c1c}{\textcolor{white}{\textbf{romance}}} & \cellcolor[HTML]{b71c1c}{\textcolor{white}{\textbf{fishing}}} & \cellcolor[HTML]{b71c1c}{\textcolor{white}{\textbf{undetermined}}}\\
\midrule
did perform & 24 & 8 & 16 & 0 & 2 & 5\\
did not perform & 0 & 1 & 4 & 0 & 0 & 1\\
indifferent & 12 & 2 & 18 & 0 & 0 & 4\\
undetermined & 0 & 1 & 2 & 0 & 0 & 0\\
\bottomrule


\caption{\label{tbl-crosstab}Table 4. Cross-tabulation of Gender
Performance by Gaming Mechanic.}

\tabularnewline
\end{longtable}

\endgroup{}

The most common cell combination were those who did perform their gender
and interacted with the land, including decorating the farm and general
farm activities (animal care, crops, rearranging, etc.). Surprisingly,
the community tab was the most spread out across gender performances,
though it was the highest with those who were indifferent to gender
performance. These players were enjoying the game in every aspect they
could, but overall gender performance favored the community aspect the
most.

\textbf{Chi-Square Result:}

A Chi-Square test of independence was conducted to examine the
relationship between Gender and Gaming Mechanic. The result was
\texttt{NA}, χ²(\texttt{13.421}) = \texttt{12}, \emph{p} =
\texttt{0.3392}. It is not statistically significant. \emph{Note: I had
to use the \texttt{chisq.test(final\$,\ final\$gaming\_mechanics)} to
find values.}

A Chi-Square test of independence was conducted to examine the
relationship between Gender and Gaming Mechanic. The result was
\texttt{NA}, χ²(\texttt{10}) = \texttt{NaN}, \emph{p} = \texttt{0.NaN}.
It is not statistically significant.

The chi\_result\$expected ran this: land mining community romance
fishing undetermined did perform 19.80 6.60 22.0 0 1.10 5.5 did not
perform 2.16 0.72 2.4 0 0.12 0.6 indifferent 12.96 4.32 14.4 0 0.72 3.6
undetermined 1.08 0.36 1.2 0 0.06 0.3

\begin{center}\rule{0.5\linewidth}{0.5pt}\end{center}

\section{Summary of Findings}\label{summary-of-findings}

\begin{itemize}
\tightlist
\item
  Community was the most preferred gaming mechanic, along with land
  closely behind.
\item
  There was an even split between those who did perform their gender and
  those who were indifferent.
\item
  There was no statistical significange between the two variables
  overall.
\end{itemize}

\begin{tcolorbox}[enhanced jigsaw, arc=.35mm, bottomrule=.15mm, bottomtitle=1mm, breakable, colback=white, colbacktitle=quarto-callout-note-color!10!white, colframe=quarto-callout-note-color-frame, coltitle=black, left=2mm, leftrule=.75mm, opacityback=0, opacitybacktitle=0.6, rightrule=.15mm, title=\textcolor{quarto-callout-note-color}{\faInfo}\hspace{0.5em}{Key Findings at a Glance}, titlerule=0mm, toprule=.15mm, toptitle=1mm]

\begin{itemize}
\item
  \textbf{RQ1} Gameplay Videos and Gender: YouTube Gameplay videos
  frames gender as indifferent to gender performance, despite their
  performance of gender or gender identity.
\item
  \textbf{RQ2} Gender Performance and Gaming Mechanic Interactions: The
  community gaming mechanic was spread out amongst all gender
  performance categories, making the comparison not statistically
  significant with p = 0.3392.
\end{itemize}

\end{tcolorbox}

\bookmarksetup{startatroot}

\chapter{Discussion \& Conclusion}\label{discussion-conclusion}

This study examined the cozy game genre through YouTube gameplay videos
of a player's first experiences. Using quantitative content analysis of
N = 100 items collected from YouTube between 04/01/2026-04/05/202, the
study first identified emergent themes through qualitative immersion
(Study 1), then applied the resulting coding scheme systematically and
tested for relationships among variables (Study 2). The findings are
discussed below in relation to each research question and the broader
theoretical framework.

\section{Addressing the Research
Questions}\label{addressing-the-research-questions}

\textbf{RQ1} asked ``How would YouTube gameplay videos frame how
different genders play cozy games?'' The qualitative analysis in Study 1
revealed that those who present as feminine were the highest part of the
dataset with those who identify as masculine were slightly under, and
the lowest being the ``other'' identifier. This included people who did
not have a preferred gender identity and/or used they/them pronouns. The
frequency distribution presented in Study 2 (Table 2, Figure 1)
confirmed that there was a fair concentration of gender performance for
those who did perform their gender and those who were indifferent to
gender performance and identity. This suggests that while those did have
a preferred gender identity, that doesn't mean they ``successfully''
performed their gender. They were split between ``did perform'' and
``indifferent,'' meaning that despite their gender identity, they
could've also leaned more towards indifferent and playing the cozy game
as it is.

This challenges Gender Performativity Theory overall. Cozy games and
their creation are supposed to be relaxing and something for an
individual to enjoy without feeling like they have to act or be a
specific way. By allowing the space to act as much as yourself as you
can, there's no inherent need to perform your gender.

\textbf{RQ2} asked ``What is the relationship between gender performance
and how the interact with their first''100 Days in Stardew Valley or
Fields of Mistria'' YouTube gameplay videos?'' The Chi-Square test
reported in Study 2 did not reveal a statistically significant
relationship between gender performance and gaming mechanics.

The absence of a significant relationship suggests that gender
performance does not vary systematically with gaming mechanics in this
dataset. This null result is itself informative: it implies that there
is no relationship between the performance of gender and the gaming
mechanics involved in a cozy game. Anything that does happen is merely
coincidental.

\section{Connecting the Two Studies}\label{connecting-the-two-studies}

The two-part design of this study allowed the qualitative and
quantitative phases to complement each other. Study 1's immersion
process ensured that the variables tested in Study 2 were grounded in
the actual patterns of the data, rather than being imposed from a
textbook or prior study. The qualitative findings held up during
immersion, as well as when it came to the Chi-Square test overall.
Individuals that were playing either \emph{Stardew Valley} or
\emph{Fields of Mistria} were playing for the sake of pure enjoyment.
They're meant to be inclusive and freeing, so during immersion, those
who were indifferent acted out of enjoyment for the game and not by
gender performance. If gender was performed, it was through verbal
interactions, but they were a miniscule part in the overall idea.

\section{Theoretical Implications}\label{theoretical-implications}

Taken together, these findings complicate the predictions of Gender
Performativity Theory. Gender Performativity Theory would predict that
repeated actions involving a person's performance of their gender would
continue through a game predominantly for women, and the present results
did not align with this prediction. This means that gender performance
is not present when it comes to the cozy game genre. Since the genre is
intended for relaxed and simpler gameplay and story, the players as
individuals can relax and enjoy themselves at a slower pace.

\section{Limitations}\label{limitations}

Several limitations of the present study should be noted. First, the
sample size was through only YouTube gameplay videos and not through
indiviual interactions of the players. Further studies could involve an
additional survey or sitting through livestreams instead of their VODs.
Second, all coding was handled by an indivudal coder. Further studies
should employ at least two-to-three as others can have different
perspectives on what exactly counts as gender identity, then report
Cohen's Kappa or Krippendorff's Alpha. Third, the study period of
04/01/2026-04/05/2026 limited how long a video could be observed for. As
this is an assignment for a class, it did not allow the overall viewing
of gameplay videos and had to be selected for the sake of time. Future
studies should allow more time in order to properly analyze gender,
Gender Performativity Theory, and the gaming mechanics.

These limitations notwithstanding, the study's systematic sampling
procedure, inductively developed codebook, and use of both qualitative
and quantitative methods represent meaningful methodological strengths.

\section{Conclusion}\label{conclusion}

The central finding of this study is that there is not a connection
between gender performance and gaming mechanics when it comes to the
cozy game genre. This matters because the genre is intended to be a
freeing and relaxing experience, which means that the pressures of
performing gender would be set aside as well. Future research might
extend this work by Future research might extend this work by studying
for a longer time to view all YouTube videos in their entirety,
additional outlets besides YouTube gameplay videos, and Uses and
Gratifications Theory by Blumler and Katz (1974) instead of Gender
Performativity Theory.

\begin{center}\rule{0.5\linewidth}{0.5pt}\end{center}

\begin{tcolorbox}[enhanced jigsaw, arc=.35mm, bottomrule=.15mm, bottomtitle=1mm, breakable, colback=white, colbacktitle=quarto-callout-note-color!10!white, colframe=quarto-callout-note-color-frame, coltitle=black, left=2mm, leftrule=.75mm, opacityback=0, opacitybacktitle=0.6, rightrule=.15mm, title=\textcolor{quarto-callout-note-color}{\faInfo}\hspace{0.5em}{A Note on This Project}, titlerule=0mm, toprule=.15mm, toptitle=1mm]

This study was conducted as a course project in MC 451: Research Methods
in Mass Media at Southern Illinois University Edwardsville (Spring
2026), under the supervision of Dr.~Alex P. Leith. The data, codebook,
and analysis files are available in the GitHub repository linked at the
top of this page.

\end{tcolorbox}

\bookmarksetup{startatroot}

\chapter{References}\label{references}

\protect\phantomsection\label{refs}
\begin{CSLReferences}{1}{0}
\bibitem[\citeproctext]{ref-GameplayVideosGoogle}
\emph{About gameplay videos - {Google Ads Help}}. (n.d.).
https://support.google.com/google-ads/answer/13268252?hl=en.

\bibitem[\citeproctext]{ref-adminGameMechanicsExplained2026}
{admin}. (2026). \emph{Game {Mechanics Explained}: {Definition}, {Types}
\& {Examples}}.

\bibitem[\citeproctext]{ref-CoregamerGenrePreferences2015}
Beyond the {``core-gamer''}: {Genre} preferences and gratifications in
computer games. (2015). \emph{Computers in Human Behavior}, \emph{44},
293--298. \url{https://doi.org/10.1016/j.chb.2014.11.020}

\bibitem[\citeproctext]{ref-boudreauWhoseVibeIt2025}
Boudreau, K., Consalvo, M., \& Phelps, A. (2025). \emph{Whose {Vibe} is
it {Anyway}? {Negotiating Definitions} of {Cozy Games}}.
\url{https://doi.org/10.24251/HICSS.2025.320}

\bibitem[\citeproctext]{ref-brtisCozyMasculinity2022}
Brtis, C. (2022). Cozy {Masculinity}. \emph{Console-Ing Passions 2022}.

\bibitem[\citeproctext]{ref-butler1988}
Butler, J. (1988). Performative acts and gender constitution: An essay
in phenomenology and feminist theory. \emph{Theatre Journal},
\emph{40}(4), 519. \url{https://doi.org/10.2307/3207893}

\bibitem[\citeproctext]{ref-coteCasualResistanceLongitudinal2020}
Cote, A. C. (2020). Casual {Resistance}: {A Longitudinal Case Study} of
{Video Gaming}'s {Gendered Construction} and {Related Audience
Perceptions}. \emph{Journal of Communication}, \emph{70}(6), 819--841.
\url{https://doi.org/10.1093/joc/jqaa028}

\bibitem[\citeproctext]{ref-fangConstructingGenderResearch2025a}
Fang, L. (2025). Constructing gender: Research on gender performativity
among video game players. \emph{Journal of Gender Studies}, 1--13.
\url{https://doi.org/10.1080/09589236.2025.2474664}

\bibitem[\citeproctext]{ref-GenderHealth}
\emph{Gender and health}. (n.d.).
https://www.who.int/health-topics/gender.

\bibitem[\citeproctext]{ref-mackayPeaceValleyMedia2025}
Mackay, O. E. L., \& Roberts, C. A. (2025). Peace in the {Valley}: {A
Media} and {Discourse Analysis} of {Eric Barone}'s {Stardew Valley
Through Utopian Theory}. \emph{Games and Culture}, \emph{20}(1), 3--19.
\url{https://doi.org/10.1177/15554120231187793}

\bibitem[\citeproctext]{ref-tangMensHarassmentBehavior2016}
Tang, W. Y., \& Fox, J. (2016). Men's harassment behavior in online
video games: {Personality} traits and game factors. \emph{Aggressive
Behavior}, \emph{42}(6), 513--521.
\url{https://doi.org/10.1002/ab.21646}

\end{CSLReferences}




\end{document}
